Despite its realism and scale, RetailBench has several limitations. First, RetailBench focuses on a single-store supermarket setting. While this setting captures many core operational pressures in retail management, it does not model multi-store coordination, competitive markets, strategic supplier behavior, labor constraints, promotions, or interactions among multiple autonomous agents. Second, although the environment is grounded in real retail product and sales signals, several dynamic components, including reviews, news events, supplier adaptation, returns, and parts of the demand model, are synthetic or rule-based. As a result, RetailBench should be viewed as a controlled stress test for long-horizon agent behavior rather than a full-fidelity simulator of real-world retail economics.
Third, our evaluation is limited to prompting-based LLM agents under a fixed environment configuration. We do not study agents with parameter updates, reinforcement learning, fine-tuning, or long-term learning across repeated episodes, all of which may lead to stronger performance. The selected-run protocol also characterizes the strongest observed behavior under the evaluated frameworks, rather than average performance across all possible agent designs or environment settings. Finally, our failure analysis is diagnostic rather than causal: the proposed metrics identify recurring bottlenecks in evidence acquisition, action conversion, product attention, and temporal follow-up, but they do not fully isolate the causal contribution of each factor. 
Addressing these limitations through richer retail environments, multi-agent extensions, learning-based adaptation, counterfactual diagnostics, and constraint-aware action control remains an important direction for future research.